\documentclass[10pt,a4paper]{amsart}
\usepackage{amsmath,amssymb,amsthm,graphicx}
% Retained minimal preamble for the Stacks Project chapter more-algebra.tex (More on Algebra).
% The chapter file pinned-original-source.tex carries only \input{preamble} on its first line, so no
% preamble text lives inside the chapter itself.  The Stacks Project uses one single shared preamble for
% every chapter; the snapshot stages it identically as sources/stacks-chow/preamble.tex and
% sources/stacks-etale-cohomology/preamble.tex.  The block below is copied mechanically, line for line,
% from sources/stacks-chow/preamble.tex lines 14-18 (the Xy-pic setup) and 154-217 (the native theorem
% declarations together with the mathematical macros).  Nothing was added, removed or redefined; the
% wrapper class and the amsmath/amssymb/amsthm/graphicx base packages are supplied by the assembler.
% Chapter: more-algebra.tex -- More on Algebra.
% For commutative diagrams we use Xy-pic
\usepackage[all]{xy}

% We use 2cell for 2-commutative diagrams.
\xyoption{2cell}

% Theorem environments.
%
\theoremstyle{plain}
\newtheorem{theorem}[subsection]{Theorem}
\newtheorem{proposition}[subsection]{Proposition}
\newtheorem{lemma}[subsection]{Lemma}

\theoremstyle{definition}
\newtheorem{definition}[subsection]{Definition}
\newtheorem{example}[subsection]{Example}
\newtheorem{exercise}[subsection]{Exercise}
\newtheorem{situation}[subsection]{Situation}

\theoremstyle{remark}
\newtheorem{remark}[subsection]{Remark}
\newtheorem{remarks}[subsection]{Remarks}

\numberwithin{equation}{subsection}

% Macros
%
\def\lim{\mathop{\mathrm{lim}}\nolimits}
\def\colim{\mathop{\mathrm{colim}}\nolimits}
\def\Spec{\mathop{\mathrm{Spec}}}
\def\Hom{\mathop{\mathrm{Hom}}\nolimits}
\def\Ext{\mathop{\mathrm{Ext}}\nolimits}
\def\SheafHom{\mathop{\mathcal{H}\!\mathit{om}}\nolimits}
\def\SheafExt{\mathop{\mathcal{E}\!\mathit{xt}}\nolimits}
\def\Sch{\mathit{Sch}}
\def\Mor{\mathop{\mathrm{Mor}}\nolimits}
\def\Ob{\mathop{\mathrm{Ob}}\nolimits}
\def\Sh{\mathop{\mathit{Sh}}\nolimits}
\def\NL{\mathop{N\!L}\nolimits}
\def\CH{\mathop{\mathrm{CH}}\nolimits}
\def\proetale{{pro\text{-}\acute{e}tale}}
\def\etale{{\acute{e}tale}}
\def\QCoh{\mathit{QCoh}}
\def\Ker{\mathop{\mathrm{Ker}}}
\def\Im{\mathop{\mathrm{Im}}}
\def\Coker{\mathop{\mathrm{Coker}}}
\def\Coim{\mathop{\mathrm{Coim}}}

% Boxtimes
%
\DeclareMathSymbol{\boxtimes}{\mathbin}{AMSa}{"02}

%
% Macros for moduli stacks/spaces
%
\def\QCohstack{\mathcal{QC}\!\mathit{oh}}
\def\Cohstack{\mathcal{C}\!\mathit{oh}}
\def\Spacesstack{\mathcal{S}\!\mathit{paces}}
\def\Quotfunctor{\mathrm{Quot}}
\def\Hilbfunctor{\mathrm{Hilb}}
\def\Curvesstack{\mathcal{C}\!\mathit{urves}}
\def\Polarizedstack{\mathcal{P}\!\mathit{olarized}}
\def\Complexesstack{\mathcal{C}\!\mathit{omplexes}}
% \Pic is the operator that assigns to X its picard group, usage \Pic(X)
% \Picardstack_{X/B} denotes the Picard stack of X over B
% \Picardfunctor_{X/B} denotes the Picard functor of X over B
\def\Pic{\mathop{\mathrm{Pic}}\nolimits}
\def\Picardstack{\mathcal{P}\!\mathit{ic}}
\def\Picardfunctor{\mathrm{Pic}}
\def\Deformationcategory{\mathcal{D}\!\mathit{ef}}

% --- appended 2026-10-12 (Chief): the project own definitions of the reference/slogan/history
% display environments, copied mechanically from the shared Stacks preamble lines 8-12
% (\usepackage{verbatim} plus the three \newenvironment lines, which the chapter preamble keeps
% verbatim from the project preamble). Nothing added, removed or redefined: these environments
% carry no mathematical content and the project itself typesets them as comments. Without these
% three lines a counted statement whose own text uses them fails with Environment slogan undefined.
\usepackage{verbatim}
\newenvironment{reference}{\comment}{\endcomment}
\newenvironment{slogan}{\comment}{\endcomment}
\newenvironment{history}{\comment}{\endcomment}
\usepackage{cleveref}
\crefname{originaldefinition}{original definition}{original definitions}
\Crefname{originaldefinition}{Original definition}{Original definitions}
\makeatletter
\newlabel{definition-pseudo-coherent}{{064Q}{source}{}{sourcepointer-1}{}}
\newlabel{definition-pseudo-coherent@cref}{{[originaldefinition][0][]064Q}{[1][0][]source}{}{}{}}
\makeatother
\crefname{originallemma}{original lemma}{original lemmas}
\Crefname{originallemma}{Original lemma}{Original lemmas}
\makeatletter
\newlabel{lemma-check-pseudo-coherent-modulo-nilpotent}{{0H76}{source}{}{sourcepointer-2}{}}
\newlabel{lemma-check-pseudo-coherent-modulo-nilpotent@cref}{{[originallemma][0][]0H76}{[1][0][]source}{}{}{}}
\makeatother
\crefname{originallemma}{original lemma}{original lemmas}
\Crefname{originallemma}{Original lemma}{Original lemmas}
\makeatletter
\newlabel{lemma-lift-complex-stably-frees}{{0BCB}{source}{}{sourcepointer-3}{}}
\newlabel{lemma-lift-complex-stably-frees@cref}{{[originallemma][0][]0BCB}{[1][0][]source}{}{}{}}
\makeatother
\crefname{originallemma}{original lemma}{original lemmas}
\Crefname{originallemma}{Original lemma}{Original lemmas}
\makeatletter
\newlabel{lemma-compute-Rlim-modules}{{091D}{source}{}{sourcepointer-4}{}}
\newlabel{lemma-compute-Rlim-modules@cref}{{[originallemma][0][]091D}{[1][0][]source}{}{}{}}
\makeatother
\crefname{originalremark}{original remark}{original remarks}
\Crefname{originalremark}{Original remark}{Original remarks}
\makeatletter
\newlabel{remark-how-unique}{{07KZ}{source}{}{sourcepointer-5}{}}
\newlabel{remark-how-unique@cref}{{[originalremark][0][]07KZ}{[1][0][]source}{}{}{}}
\makeatother
\title{A derived limit of a compatible system of pseudo-coherent objects over surjective ring maps is pseudo-coherent}
\author{The Stacks Project Authors}
\date{stacks-project:more-algebra.tex:local-snapshot-20261012; GNU FDL 1.2 or later}
\begin{document}
\maketitle
\setcounter{subsection}{1}
\noindent\textit{Editorial note.} All delivered mathematical text is the author's own, copied from the pinned Stacks Project chapter file more-algebra.tex; no statement or proof step was rewritten, reordered or removed. The counted statement is the labeled lemma environment at source lines 28335-28351 (canonical tag 0CQF) and the counted proof is the proof environment at source lines 28353-28375, which immediately follows it and is its own proof.\par\medskip
\noindent\textit{Editorial note.} The five references are four earlier results of this chapter plus a remark of this chapter cited alongside one of them; the kinds recorded are the type words appearing in the labels themselves, following the lane convention, and none of them is the name of the environment in which the label lives.\par\medskip
\noindent\textit{Editorial note.} The preamble span runs to line 90, the current end of the retained preamble, which includes the appended environment definitions; the counted text does not use them, and the span is taken whole so that the file version is unambiguous.\par\medskip
\noindent\textit{Editorial note.} The source URL is the source-level Stacks Project URL, as required when more than one card is produced from one source; the per-claim tag is carried by claim\_key and by this note. The counted statement and proof cite no bibliography entry, so this card needs no bibliography.\par\medskip
\noindent\textit{Allowed standard prerequisites (editorial).}\begin{itemize}
\item A pseudo-coherent object over a ring is represented by a bounded above complex of finite free modules, and pseudo-coherence is preserved under a surjection with locally nilpotent kernel when the target object is pseudo-coherent. invoked at proof lines 28353-28359. Produces the representative over the first ring and propagates pseudo-coherence to every member of the system.
\item Maps in the derived category between the members of the system may be lifted to maps of complexes of finite free modules inducing isomorphisms after reduction to the previous ring. invoked at proof lines 28359-28365. Builds the system of complexes whose limit represents the derived limit.
\item The derived limit of a system of complexes is represented by the limit of the complexes, and the limit of finite free modules of bounded rank is again finite free of that rank. used at proof lines 28365-28374. Shows that the limit complex is a bounded above complex of finite free modules, which gives pseudo-coherence and the final assertion.
\end{itemize}\medskip
\ifdefined\hypertarget\hypertarget{sourcepointer-1}{}\fi
\noindent\textit{Source reference.} Original definition 064Q: sources/stacks-more-algebra/pinned-original-source.tex, source lines 15874--15874. The reference supplies the definition used at proof lines 28353-28355 to choose the bounded above complex of finite free modules; the induction built on it is inside the counted span.\par\medskip
\ifdefined\hypertarget\hypertarget{sourcepointer-2}{}\fi
\noindent\textit{Source reference.} Original lemma 0H76: sources/stacks-more-algebra/pinned-original-source.tex, source lines 19498--19498. The reference is applied at proof lines 28355-28357 to conclude that every object of the system is pseudo-coherent; the lifting argument that follows is inside the counted span.\par\medskip
\ifdefined\hypertarget\hypertarget{sourcepointer-3}{}\fi
\noindent\textit{Source reference.} Original lemma 0BCB: sources/stacks-more-algebra/pinned-original-source.tex, source lines 19517--19517. The reference is invoked at proof lines 28359-28361 to construct the complexes and maps of the system with strict reductions; the limit computation that uses them is inside the counted span.\par\medskip
\ifdefined\hypertarget\hypertarget{sourcepointer-4}{}\fi
\noindent\textit{Source reference.} Original lemma 091D: sources/stacks-more-algebra/pinned-original-source.tex, source lines 23574--23574. The reference identifies the derived limit with the limit of the system of complexes, used at proof lines 28365-28367; the reading off of finite freeness afterwards is inside the counted span.\par\medskip
\ifdefined\hypertarget\hypertarget{sourcepointer-5}{}\fi
\noindent\textit{Source reference.} Original remark 07KZ: sources/stacks-more-algebra/pinned-original-source.tex, source lines 23753--23753. The remark is cited jointly with the previous reference at proof lines 28365-28367 and carries no separate mathematical input; it is a pointer accompanying that identification.\par\medskip
\begin{lemma}
\label{lemma-Rlim-pseudo-coherent-gives-pseudo-coherent}
Let $A = \lim A_n$ be a limit of an inverse system $(A_n)$ of rings.
Suppose given $K_n \in D(A_n)$ and maps $K_{n + 1} \to K_n$
in $D(A_{n + 1})$. Assume
\begin{enumerate}
\item the transition maps $A_{n + 1} \to A_n$ are surjective
with locally nilpotent kernels,
\item either all $K_n$ are pseudo-coherent or there exists an integer
$n_0$ such that $K_{n_0}$ is pseudo-coherent and the kernels of
$A_{n + 1} \to A_n$ are nilpotent ideals for $n \geq n_0$,
\item the maps induce isomorphisms
$K_{n + 1} \otimes_{A_{n + 1}}^\mathbf{L} A_n \to K_n$.
\end{enumerate}
Then $K = R\lim K_n$ is a pseudo-coherent object of $D(A)$
and $K \otimes_A^\mathbf{L} A_n \to K_n$ is an isomorphism for all $n$.
\end{lemma}

\begin{proof}
By assumption we can find a bounded above complex of
finite free $A_1$-modules $P_1^\bullet$ representing $K_1$, see
Definition \ref{definition-pseudo-coherent}. By
Lemma \ref{lemma-check-pseudo-coherent-modulo-nilpotent}
we conclude that $K_n$ is pseudo-coherent for all $n$.
Then by Lemma \ref{lemma-lift-complex-stably-frees}
we can, by induction on $n > 1$, find
complexes $P_n^\bullet$ of finite free $A_n$-modules representing $K_n$
and maps $P_n^\bullet \to P_{n - 1}^\bullet$ representing the maps
$K_n \to K_{n - 1}$ inducing isomorphisms (!)
of complexes $P_n^\bullet \otimes_{A_n} A_{n - 1} \to P_{n - 1}^\bullet$.
Thus $K = R\lim K_n$ is represented by $P^\bullet = \lim P_n^\bullet$, see
Lemma \ref{lemma-compute-Rlim-modules} and
Remark \ref{remark-how-unique}.
Since $P_n^i$ is a finite free $A_n$-module for each $n$ and
$A = \lim A_n$ we see that $P^i$ is finite free of the same rank
as $P_1^i$ for each $i$.
This means that $K$ is pseudo-coherent.
It also follows that $K \otimes_A^\mathbf{L} A_n$ is represented by
$P^\bullet \otimes_A A_n = P_n^\bullet$ which proves the final
assertion.
\end{proof}

\end{document}
